\documentclass[11pt]{article}

\usepackage[T1]{fontenc}
\usepackage[utf8]{inputenc}
\usepackage[english]{babel}
\usepackage{amsmath,amssymb,amsthm}
\usepackage{booktabs}
\usepackage{geometry}
\usepackage{hyperref}
\usepackage{url}

\geometry{margin=1in}

\hypersetup{
  colorlinks=true,
  linkcolor=black,
  citecolor=blue,
  urlcolor=blue,
  pdftitle={Computational Extension of OEIS Sequence A224540: New Terms Through a(28) via a Truncated Inverse Collatz Tree},
  pdfauthor={Carlo Corti}
}

\newtheorem{theorem}{Theorem}
\newtheorem{proposition}[theorem]{Proposition}
\newtheorem{lemma}[theorem]{Lemma}
\newtheorem{corollary}[theorem]{Corollary}
\theoremstyle{definition}
\newtheorem{definition}[theorem]{Definition}
\theoremstyle{remark}
\newtheorem{remark}[theorem]{Remark}

\title{Computational Extension of OEIS Sequence A224540:\\
New Terms Through \(a(28)\) via a Truncated Inverse Collatz Tree}
\author{Carlo Corti\\
Independent researcher\\
\texttt{carlo.corti@outlook.com}}
\date{June 2026}

\begin{document}

\maketitle

\begin{abstract}
OEIS A224540 counts the positive integers \(k\) whose Collatz trajectory,
stopped at the first occurrence of \(1\), never exceeds \(3^n\).  The public
OEIS b-file previously contained terms through \(a(21)\).  This paper reports
a computational extension through \(a(28)\), with new values
\[
\begin{aligned}
a(22)&=12406199367, & a(23)&=37216010685,\\
a(24)&=111651745638, & a(25)&=334951946950,\\
a(26)&=1004855123675, & a(27)&=3014552251123,\\
a(28)&=9043690122807.
\end{aligned}
\]
The computation used an iterative depth-first traversal of the inverse
Collatz tree rooted at the terminal value \(1\), truncated by the inclusive
bound \(M=3^n\).  The largest terms were produced by segmented subtree
campaigns and checked by an independent artifact-level validator that
reparsed the result files, recomputed powers \(3^n\), verified segment
coverage, and recomputed aggregate sums.  The independent validation was not
a second full traversal of all inverse subtrees for \(n=28\).
\end{abstract}

\noindent\textbf{2020 Mathematics Subject Classification.}
Primary 11Y55; Secondary 11B83, 11Y35.

\noindent\textbf{Key words and phrases.}
Collatz iteration, \(3x+1\) problem, integer sequence, OEIS, inverse tree,
computational number theory.

\section{Introduction}

Let \(T\) denote the usual Collatz map,
\[
T(k)=
\begin{cases}
k/2, & k \equiv 0 \pmod 2,\\
3k+1, & k \equiv 1 \pmod 2.
\end{cases}
\]
For a starting value \(k\) whose forward orbit reaches \(1\), the stopped
trajectory is the finite sequence \(k,T(k),T^2(k),\ldots,1\), truncated at the
first occurrence of \(1\).
OEIS A224540 is the sequence
\[
a(n)=\#\{k\geq 1:\hbox{ the trajectory of }k\hbox{, stopped at the first }1,
\hbox{ never exceeds }3^n\}.
\]
Equivalently, \(a(n)\) counts the positive starting values \(k\) whose stopped
trajectory maximum is at most \(3^n\).  The inequality is inclusive.
For \(k=1\), the stopped trajectory is the one-term trajectory \([1]\), so
\(1\) is counted for every \(n\geq 0\).

For example, for \(n=3\), the bound is \(3^3=27\), and the twelve values listed
in the OEIS entry are
\[
1,2,3,4,5,6,8,10,12,16,20,24.
\]
The local OEIS snapshot used for this project had offset \(0,2\), listed terms
through \(a(21)\), and credited T. D. Noe as the author of the entry, with
\(a(20)\) and \(a(21)\) supplied by Donovan Johnson \cite{oeisA224540}.

The direct definition suggests scanning all \(k\leq 3^n\) and computing
trajectory maxima.  The production computation instead enumerates exactly the
accepted values by traversing the inverse Collatz tree from the terminal root
\(1\), with the same ceiling \(M=3^n\).  This gives a simple finite graph
problem for each \(n\), avoids a visited set, and keeps memory usage governed
by an explicit DFS stack rather than by the size of the accepted set.

The contribution of this paper is the extension of the b-file through
\(n=28\), together with the finite computation and validation artifacts that
support it.  The result is computational: it does not assert anything about
the Collatz conjecture beyond the finite stopped trajectories represented in
the truncated inverse tree.  Standard background on the \(3x+1\) problem is
given by Lagarias \cite{lagarias1985,lagarias2010}.

\section{Definition and inverse-tree framework}

\begin{definition}
For \(M\geq 1\), let \(S_M\) be the smallest set of positive integers obtained
as follows.  The root \(1\) belongs to \(S_M\).  If \(x\in S_M\), then the
following inverse predecessors are admitted when they satisfy the displayed
conditions:
\[
2x \quad\hbox{if } 2x\leq M,
\]
and
\[
\frac{x-1}{3} \quad\hbox{if } x\equiv 1\pmod 3,\quad
\frac{x-1}{3}\hbox{ is odd},\quad
\frac{x-1}{3}>1,\quad
\frac{x-1}{3}\leq M.
\]
\end{definition}

\begin{remark}
The condition \((x-1)/3>1\) is part of the stopped-at-\(1\) convention.  Without
it, the odd predecessor \(1\) of \(4\) would create the trivial cycle
\(1,4,2,1\) inside the inverse construction.  In A224540 the trajectory is
stopped when \(1\) is first reached, so \(1\) is a terminal root, not a node
whose outgoing Collatz step is followed again.
\end{remark}

\begin{proposition}[Truncated inverse-tree equivalence]
\label{prop:inverse-tree}
For every \(M\geq 1\), \(S_M\) is exactly the set of positive integers
\(k\leq M\) whose Collatz trajectory, stopped at the first occurrence of \(1\),
never exceeds \(M\).
\end{proposition}

\begin{proof}
If \(y\) is admitted as an inverse predecessor of \(x\), then \(T(y)=x\):
this is immediate for the doubled predecessor \(y=2x\), and for the odd
predecessor \(y=(x-1)/3\) it follows from the imposed congruence and parity
conditions.  Therefore every node generated from the root has a finite forward
path to \(1\), obtained by reversing its generation path.  Every node on that
path was admitted only after the bound check \(y\leq M\).  Hence every
generated node satisfies the defining property.

Conversely, let \(k\leq M\) have a stopped trajectory
\[
k, T(k), T^2(k),\ldots,1
\]
all of whose entries are at most \(M\).  Reversing this finite path starts at
\(1\).  Each reversed step is either \(x\mapsto 2x\), for an even predecessor,
or \(x\mapsto (x-1)/3\), for an odd predecessor.  The predecessor is odd in the
second case; moreover, if that predecessor is \(y\), then \(x=T(y)=3y+1\), so
\(x\equiv 1\pmod 3\).  It remains at most \(M\) by hypothesis and is greater
than \(1\) because the value \(1\) occurs only as the terminal element of the
stopped trajectory.  Thus every element of the stopped trajectory, in
particular \(k\), is generated in \(S_M\).

The next Collatz iterate of any positive integer is unique.  Therefore a
generated node \(y\neq 1\) cannot enter the rooted inverse construction under
two different parents \(x\), where the parent is the forward image \(T(y)=x\)
in the rooted tree.  The rule excluding the odd predecessor \(1\) of \(4\)
removes the only terminal-cycle re-entry relevant to the stopped convention.
Consequently the inverse traversal counts each accepted starting value exactly
once.
\end{proof}

Taking \(M=3^n\), Proposition~\ref{prop:inverse-tree} gives
\[
a(n)=\#S_{3^n}.
\]
This is the mathematical basis for the production program.

\section{Production traversal and segmentation}

The production source is a C17 program, stored as \path{src/a224540.c} in the
public GitHub/Zenodo package \cite{cortiA224540Package2026}.  For a fixed
\(M=3^n\), it performs an iterative DFS over the inverse tree:
\[
\begin{array}{l}
\hbox{initialize stack with }1;\\
\hbox{while the stack is nonempty, pop }x\hbox{ and count it;}\\
\hbox{push }2x\hbox{ if }2x\leq M;\\
\hbox{push }(x-1)/3\hbox{ under the odd-predecessor conditions above.}
\end{array}
\]
The implementation forms the doubling test as \(x\leq M/2\), so that \(2x\)
is not computed before the bound check.  It stores states and counts in
\texttt{uint64\_t} within the range used here, and uses \texttt{\_\_uint128\_t}
for decimal parsing, powers of \(3\), and overflow checks.  The project
numeric ledger records that \(3^{28}=22876792454961\), far below the
\texttt{uint64\_t} limit; the code rejects powers that do not fit.

The program has several modes relevant to reproducibility.  A replay mode
recomputes the inverse-tree counts for the compiled known prefix
\(n=0,\ldots,21\) and compares them with the OEIS-prefix rows in the b-file.
A range mode produced \(a(22)\) through \(a(25)\).  For the larger terms, the
traversal was split at a fixed depth into a finite frontier, and contiguous
intervals of frontier indices were processed as independent segment jobs.
Aggregation then added the prefix count and the segment subtree counts, after
checking that segment intervals covered the whole frontier without gaps or
overlaps.  Table~\ref{tab:campaign-summary} summarizes the production
campaigns.

For segmentation, the depth of the root \(1\) is \(0\), and the depth of a
node is the number of inverse steps from the root.  For a chosen depth \(d\),
the frontier is the ordered list of nodes first encountered at exactly depth
\(d\) during the same deterministic iterative DFS used by the program.  The
child-generation order is the doubled child first, then the odd predecessor
when it exists; the stack is last-in, first-out, so when both children exist
the odd predecessor is visited before the doubled child.  The frontier order
is this pop/visit order, not merely the textual order of the push operations.
The prefix count is the number of nodes at depths \(<d\).  Frontier indices are
zero-based.  A segment with interval \([i,j]\) processes the complete subtrees
rooted at frontier indices \(i,i+1,\ldots,j\), inclusive.  Since
Proposition~\ref{prop:inverse-tree} gives a rooted tree, these segment subtrees
are disjoint, and the union of all frontier intervals together with the prefix
nodes is the full truncated tree.

Segmentation was used for resumability and auditability, not as a load
balancing claim.  For \(n=28\), the frontier at depth \(24\) contained only
179 roots, but the subtrees rooted at those frontier nodes had very different
sizes.  The initial equal-size plan was refined by splitting one adjacent
frontier region, so the final aggregate uses fourteen segment files; the
complete ledger is Table~\ref{tab:n28-segments}.  Two segment files account
for most of the wall-clock elapsed time: frontier indices \(0\)--\(14\)
recorded 4021158357110 counted nodes in 11951.82 seconds, while indices
\(45\)--\(59\) recorded 4064690194873 counted nodes in 12149.51 seconds.
Other segments with similar numbers of roots were much smaller; for example,
indices \(165\)--\(178\) recorded 4732930237 nodes in 14.27 seconds.  This
imbalance is an operational property of the truncated inverse Collatz tree, not
an aggregation error.

The production environment for the recorded local campaign was an AMD Ryzen 9
7940HS system with 16 logical CPUs and 64 GB DDR5 installed RAM, running
Linux Mint 22.3 on x86-64.  The underlying distribution base reported by
kernel or package tooling is not used as the operating-system description for
this paper.  The current package builds with GCC 13.3.0 as a C17 program; the
release Makefile uses the optimization flags \texttt{-O3},
\texttt{-march=native}, \texttt{-flto}, \texttt{-funroll-loops},
\texttt{-fomit-frame-pointer}, and \texttt{-DNDEBUG}, plus warning flags.
For checked intermediate arithmetic, the implementation uses the GCC/Clang
extension \texttt{\_\_uint128\_t} under this documented GCC build environment.
The validation target recorded in the package was
\[
\texttt{make -C src all validate sanitize RUN\_ID=20260630Tbfile\_extended\_final}.
\]
Here \texttt{validate} is the full validation target; in the public Makefile,
\texttt{test} is an alias for \texttt{validate}.  It runs the production-program
self-tests, CLI negative tests, known-prefix replay, calibration, small
segmented-versus-monolithic checks, aggregate adversarial tests, the independent
top-down check through \(n=14\), and local checksum generation.  The separate
artifact-level validator is the Python script described below.

\begin{table}[ht]
\centering
\begin{tabular}{rrrrrr}
\toprule
\(n\) & \(M=3^n\) & mode & depth & segments & elapsed seconds \\
\midrule
22--25 & \(3^{22}\)--\(3^{25}\) & range & -- & 0 & 1634.34 \\
26 & 2541865828329 & segmented & 18 & 1 & 3021.58 \\
27 & 7625597484987 & segmented & 20 & 2 & 9039.15 \\
28 & 22876792454961 & segmented & 24 & 14 & 26987.87 \\
\bottomrule
\end{tabular}
\caption{Production campaign summary from the run manifest TSV, with elapsed
times rounded to two decimals.  The first row is a single range campaign for
\(n=22,\ldots,25\).  Elapsed times are wall-clock measurements from the run
manifests; rounded segment displays need not sum exactly to a separately
recorded rounded campaign total.}
\label{tab:campaign-summary}
\end{table}

\section{Validation and certification boundary}

The validation record has three layers.

First, the production program contains local child-rule tests, branch-order
invariance tests, and segmented-versus-monolithic checks on small instances.
Here branch-order invariance means that small counts agree under the supported
child-order variants, while segmented-versus-monolithic checks compare a small
frontier partition with a single unsegmented count.  The Makefile target
\texttt{make -C src test}, an alias of \texttt{make -C src validate}, runs CLI
tests, replay of the compiled known prefix, a small independent top-down check
through \(n=14\), segmented/monolithic checks at \(n=14\) and depth \(8\), and
aggregate adversarial checks.  The top-down check is a separate Python
enumerator used only at small scale.  The target \texttt{make -C src sanitize}
builds and runs an
AddressSanitizer/UndefinedBehaviorSanitizer smoke validation.

Second, the known OEIS prefix \(a(0),\ldots,a(21)\) was replayed.  In the
published raw package, \path{data/certified_terms.tsv} records these terms as
coming from the OEIS b-file replay oracle, while \(a(22),\ldots,a(25)\) are
from the range campaign and \(a(26),a(27),a(28)\) are from segmented aggregate
campaigns.

Third, the script \path{tools/a224540_validate_new_results.py} performs an
independent artifact-level validation.  It does not call the production
binary.  It parses the range and aggregate TSV files, validates that
\path{data/b224540.txt} is contiguous for \(n=0,\ldots,28\), recomputes the
powers \(3^n\), reconstructs the finite frontier summaries for segmented
campaigns, checks complete segment coverage, and recomputes aggregate sums.
For segmented campaigns, complete coverage means that the recorded frontier
intervals partition the zero-based frontier without gaps or overlaps; the
validator also checks that the aggregate total is the prefix count plus the
recorded segment subtree counts.  Its recorded verdict is PASS.  The status
\texttt{AGGREGATED\_NO\_ORACLE} denotes an aggregate reconstructed from segment
artifacts without a separate external oracle for that large term.

\begin{table}[ht]
\centering
\begin{tabular}{rrrrrr}
\toprule
\(n\) & depth & prefix & frontier & segments & total \\
\midrule
26 & 18 & 168 & 44 & 1 & 1004855123675 \\
27 & 20 & 270 & 72 & 2 & 3014552251123 \\
28 & 24 & 689 & 179 & 14 & 9043690122807 \\
\bottomrule
\end{tabular}
\caption{Segmented campaign checks recomputed by the independent artifact-level
validator.  The totals include the prefix count and all segment subtree counts.}
\label{tab:segment-checks}
\end{table}

\begin{table}[ht]
\centering
\begin{tabular}{rrrr}
\toprule
frontier interval & roots & subtree count & elapsed seconds \\
\midrule
0--14 & 15 & 4021158357110 & 11951.82 \\
15--15 & 1 & 9437625941 & 28.28 \\
16--20 & 5 & 113738272773 & 345.38 \\
21--29 & 9 & 30709962147 & 93.37 \\
30--44 & 15 & 133937400225 & 406.50 \\
45--59 & 15 & 4064690194873 & 12149.51 \\
60--74 & 15 & 89293833363 & 265.38 \\
75--89 & 15 & 19101684797 & 57.79 \\
90--104 & 15 & 179424181121 & 537.10 \\
105--119 & 15 & 32810854681 & 98.37 \\
120--134 & 15 & 179352408944 & 543.57 \\
135--149 & 15 & 157417194040 & 472.74 \\
150--164 & 15 & 7885221866 & 23.78 \\
165--178 & 14 & 4732930237 & 14.27 \\
\bottomrule
\end{tabular}
\caption{Complete \(n=28\) segment ledger from the campaign TSVs, with elapsed
times rounded to two decimals for readability.  Together with the prefix count
\(689\), the subtree counts sum to \(a(28)=9043690122807\).}
\label{tab:n28-segments}
\end{table}

The exhaustiveness of the finite enumeration for each \(n\) rests on the
production kernel and the recorded campaign manifests: every node of the
corresponding truncated inverse tree is counted once by the traversal or by a
complete frontier partition.  The independent validator certifies the
consistency of the produced artifacts and the segment aggregation, but it is
not a second production-scale traversal of all inverse subtrees for \(n=28\).
This distinction is part of the certification boundary.

More explicitly, the independent validator checks file syntax, the contiguous
b-file through \(n=28\), recomputed powers \(3^n\), consistency of the range
and aggregate TSVs, complete segment coverage, and aggregate sums.  It does
not prove the absence of a systematic defect in the production traversal by
rerunning the full computation with a distinct implementation.  The supporting
evidence for that part is the reviewed C17 kernel, replay of the known OEIS
prefix, small independent top-down checks, sanitizer runs, segmented versus
monolithic tests on small instances, adversarial aggregate tests, and the
recorded production manifests.

\section{Results}

\begin{theorem}[Artifact-checked computational extension through \(a(28)\)]
\label{thm:main}
Subject to the correct execution of the reviewed production kernel and the
recorded campaign manifests in the raw package, OEIS A224540 has the following
values for \(n=22,\ldots,28\):
\[
\begin{array}{c|r}
n & a(n)\\
\hline
22 & 12406199367\\
23 & 37216010685\\
24 & 111651745638\\
25 & 334951946950\\
26 & 1004855123675\\
27 & 3014552251123\\
28 & 9043690122807
\end{array}
\]
Together with the replayed OEIS prefix, the b-file is extended contiguously
through \(n=28\).
\end{theorem}

\begin{proof}
For each \(n\), Proposition~\ref{prop:inverse-tree} reduces \(a(n)\) to the
finite count of the truncated inverse tree \(S_{3^n}\), with the inclusive
bound \(x\leq 3^n\) and the terminal-root convention at \(1\).  The production
program implements exactly these inverse-child rules and counts the generated
nodes by DFS.

For \(n=22,\ldots,25\), the range campaign TSV
\path{validation/campaigns/a22_a25/a22_a25.tsv} records the four counts in the
theorem with status \texttt{COMPUTED}.  For \(n=26,27,28\), the aggregate TSVs
under \path{validation/campaigns/} record the segmented totals in
Table~\ref{tab:segment-checks} with status \texttt{AGGREGATED\_NO\_ORACLE}.
The aggregate records contain complete frontier coverage data, and the
independent validator recomputed the corresponding powers, prefixes,
frontiers, segment coverage, and sums.  The same validator checked that the
extended b-file agrees with the result table for all \(n=0,\ldots,28\).
Therefore the recorded finite computations give the displayed extension.
\end{proof}

\begin{corollary}
The OEIS-style b-file \path{data/b224540.txt} in the raw package is a
contiguous table of \(n,a(n)\) for \(0\leq n\leq 28\), with the seven new terms
listed in Theorem~\ref{thm:main}.  Table~\ref{tab:terms} summarizes the
certification source for each part of the extended table.
\end{corollary}

\begin{table}[ht]
\centering
\begin{tabular}{rrl}
\toprule
\(n\) & \(a(n)\) & source in package \\
\midrule
0--21 & OEIS prefix through 4135278456 & b-file replay oracle \\
22 & 12406199367 & range campaign \\
23 & 37216010685 & range campaign \\
24 & 111651745638 & range campaign \\
25 & 334951946950 & range campaign \\
26 & 1004855123675 & segmented aggregate \\
27 & 3014552251123 & segmented aggregate \\
28 & 9043690122807 & segmented aggregate \\
\bottomrule
\end{tabular}
\caption{Certification sources summarized from the certified-terms and result
TSV files.}
\label{tab:terms}
\end{table}

\section{Conclusion}

The computation extends A224540 from the previously recorded \(a(21)\) to
\(a(28)\).  The main mathematical point is that the stopped-at-\(1\) Collatz
condition is exactly represented by a finite inverse tree rooted at terminal
\(1\), provided the predecessor \(1\) of \(4\) is excluded and the bound
\(x\leq 3^n\) is treated inclusively.  The largest reported term,
\[
a(28)=9043690122807
\]
comes from a segmented traversal of the truncated tree at
\(M=3^{28}=22876792454961\).  The validation artifacts support the b-file
extension through \(n=28\), with the explicit limitation that the independent
validator checks artifact consistency and aggregation rather than repeating a
second full production-scale traversal.

\section*{Data availability}

The code, b-file, result tables, campaign TSVs, manifests, validation
summaries, and independent artifact-level validator are part of the public
GitHub repository and Zenodo archive listed in the References.
The essential reproducibility artifacts are \path{src/a224540.c},
\path{src/Makefile}, \path{data/b224540.txt},
\path{data/certified_terms.tsv}, \path{results/A224540_terms.tsv},
\path{validation/run_manifest.tsv}, the \path{validation/campaigns/}
subtree, and \path{tools/a224540_validate_new_results.py}.  The canonical
OEIS-style b-file in the package is \path{data/b224540.txt}; result tables and
certification ledgers are cross-checking artifacts, not competing b-files.

The GitHub URL and Zenodo DOI in the reference below identify the public
release package.

\section*{AI use disclosure}

Large language models were used as research and drafting assistants during the
A224540 workflow, including OpenAI Codex and external LLM-based review systems
from the ChatGPT, Claude, DeepSeek, Gemini, Grok, Kimi, Qwen, MiniMax and Z
families.  They are named at family level where exact model-version metadata
were not consistently available in the local review artifacts.

Mathematical analysis.  LLM assistance helped organize the stopped-trajectory
convention, the inverse-tree equivalence, the terminal-root edge case, and the
certification boundary.  The mathematical claims used in this manuscript were
checked against the production source, the OEIS snapshot, and the validation
artifacts.

Algorithm and code.  LLM tools assisted with code review, validation planning,
and manuscript-facing explanations of the C17 implementation.  The reported
values do not depend on prose generated by an AI system; they are taken from
the recorded production TSVs, b-file, manifest, and independent artifact-level
validator.

Manuscript.  LLM assistance was used while drafting and reviewing this TeX
manuscript.  Numerical statements in the manuscript were cross-checked against
the local package files before inclusion.

The author takes full scientific responsibility for the mathematical
statements, algorithm, implementation, numerical results, validation
boundaries, and conclusions of this paper.

\begin{thebibliography}{9}
\sloppy
\raggedright

\bibitem{oeisA224540}
OEIS Foundation Inc.,
\emph{Sequence A224540: Number of numbers \(k\) such that all terms of the
Collatz \((3x+1)\) iteration of \(k\) are \(\leq 3^n\)},
The On-Line Encyclopedia of Integer Sequences,
\url{https://oeis.org/A224540}.
Local project snapshot accessed June 30, 2026 listed offset \(0,2\), terms
through \(a(21)\), author T. D. Noe, Apr. 24, 2013, and extension
\(a(20)\)--\(a(21)\) by Donovan Johnson, Jun. 5, 2013.

\bibitem{lagarias1985}
J. C. Lagarias,
\emph{The \(3x+1\) problem and its generalizations},
American Mathematical Monthly \textbf{92} (1985), no.~1, 3--23.
\url{https://doi.org/10.1080/00029890.1985.11971528}.

\bibitem{lagarias2010}
J. C. Lagarias, editor,
\emph{The Ultimate Challenge: The \(3x+1\) Problem},
American Mathematical Society, Providence, RI, 2010.
ISBN-13: 9780821849408; ISBN-10: 0821849409.

\bibitem{cortiA224540Package2026}
C. Corti,
\emph{A224540: Computational extension through \(a(28)\) by truncated inverse
Collatz-tree traversal},
public GitHub repository and Zenodo archive, 2026.
\url{https://github.com/carcorti/A224540}.
\url{https://doi.org/10.5281/zenodo.21096334}.

\end{thebibliography}

\end{document}
